\documentclass[12pt,aspectratio=169]{beamer}
\input{header_beamer}
\usetheme{Boadilla}
\usecolortheme{beaver}
\setbeamertemplate{page number in head/foot}{}
\setbeamertemplate{navigation symbols}{}

\title{Lecture-07: Random Processes}
\date{}

\begin{document}
\pagestyle{empty}
\maketitle

\begin{frame}{Introduction} 
\begin{block}{Reminder}<1->
For an arbitrary index set $T$, and a real-valued function $x \in \R^T$, 
the projection operator $\pi_t:\R^T\to\R$ maps $x\in\R^T$ to $\pi_t(x) = x_t$. 
\end{block}

\begin{Definition}[Random process]<2->
Let $(\Omega, \sF, P)$ be a probability space. 
For an arbitrary index set $T$ and state space $\sX \subseteq \R$, 
a map $X : \Omega \to \sX^T$ is called a \textbf{random process} if the projections $X_t:\Omega\to\sX$ defined by $\omega \mapsto X_t(\omega) \triangleq (\pi_t\circ X)(\omega)$ are random variables on the given probability space. 
\end{Definition}
\end{frame}

\begin{frame}{Random Processes}
\begin{Definition} [Sample Path]<1->
For each outcome $\omega \in \Omega$, we have a function $X(\omega): T \mapsto \sX$ called the \textbf{sample path} or the \textbf{sample function} of the process $X$. %, 
\end{Definition}
\begin{block}{Reminder}<2-> 
A random process $X$ defined on probability space $(\Omega, \sF, P)$ with index set $T$ and state space $\sX\subseteq\R$, can be thought of as 
\begin{enumerate}[(a)]
\item <3-> a map $X: \Omega \times T \to \sX$, 
\item <4-> a map $X: T \to \sX^\Omega$, i.e. a collection of random variables $X_t: \Omega \to \sX$ for each time $t \in T$, 
\item <5-> a map $X: \Omega \to \sX^T$, i.e. a collection of sample functions $X(\omega): T \to \sX$ for each random outcome $\omega \in \Omega$. 
\end{enumerate}
\end{block}
\end{frame}

\begin{frame}
\begin{block}{Classification}
\begin{itemize}
\item <1-> State space $\sX$ can be countable or uncountable, corresponding to discrete or continuous valued process.  
\item <2-> If the index set $T \subseteq \R$ is countable, the stochastic process is called \textbf{discrete}-time stochastic process or random sequence. 
\item <3-> When the index set $T$ is uncountable, it is called \textbf{continuous}-time stochastic process. 
\item <4-> The index set $T$ doesn't have to be time, if the index set is space, and then the stochastic process is spatial process. 
\item <5-> When $T = \R^n \times [0, \infty)$, stochastic process $X$ is a spatio-temporal process. 
\end{itemize}
\end{block}
\end{frame}

\begin{frame}{Random Processes}
\begin{Example} 
\begin{enumerate}[i\_]
\item <1-> Discrete random sequence: brand switching, discrete time queues, number of people at bank each day.
\item <2-> Continuous random sequence: stock prices, currency exchange rates, waiting time in queue of $n$th arrival, workload at arrivals in time sharing computer systems.
\item <3-> Discrete random process:  counting processes, population sampled at birth-death instants, number of people in queues.
\item <4-> Continuous random process: water level in a dam, waiting time till service in a queue, location of a mobile node in a network.
\end{enumerate}
\end{Example}
\end{frame}

\begin{frame}{Measurability}
\begin{block}{Note}
\begin{itemize}
\item <1-> For random process $X:\Omega\to \sX^T$ defined on  $(\Omega, \sF, P)$, the projections $X_t \triangleq \pi_t\circ X$ are $\sF$-measurable random variables. Therefore, $A_{X_t}(x) \triangleq X_t^{-1}(-\infty, x] \in \sF$ for all $t \in T$ and $x \in \R$. \end{itemize}
\end{block}
\begin{Definition}<2->[Measurability of a random process] 
A random map $X:\Omega\to\sX^T$ is called $\sF$-\textbf{measurable},  
if $A_{X_t}(x) = X_t^{-1}(-\infty, x]\in \sF$ for all $t \in T$ and $x\in \R$. 
\end{Definition} 
\end{frame}

\begin{frame}{Measurability}
\begin{Definition}<1-> 
The \textbf{event space generated by a random process} $X: \Omega\to\sX^T$ defined on a probability space $(\Omega,\sF,P)$ is given by 
\EQ{
\sigma(X) \triangleq \sigma(A_{X_t}(x): t\in T, x \in \R). 
}
\end{Definition}
\begin{Definition}<2->
For a random process $X:\Omega \to \sX^T$ defined on the probability space $(\Omega,\sF,P)$,  
we define the projection of $X$ onto components $S \subseteq T$ as the random vector $X_S: \Omega \to \sX^S$, where $X_S \triangleq (X_s: s \in S)$.
\end{Definition}
\end{frame}

\begin{frame}{Measurability}
\begin{block}{Reminder}
\begin{itemize}
\item <1-> Recall that $\pi_t^{-1}(-\infty, x] = \bigtimes_{s \in T}(-\infty, x_s]$ where $x_s = x$ for $s = t$ and $x_s = \infty$ for all $s\neq t$. 
\item <2-> The $\sF$-measurability of process $X$ implies that for any countable set $S \subseteq T$, 
we have $A_{X_S}(x_S) \triangleq \cap_{s \in S}A_{X_s}(x_s) \in \sF$ for $x_S \in \sX^S$. 
\item <3-> We can define $A_X(x) \triangleq \cap_{t \in T}A_{X_t}(x_t)$ for any $x \in \R^T$. 
\item <4-> However, $A_X(x)$ is guaranteed to be an event only when $S \triangleq \set{t \in T: \pi_t(x) < \infty}$ is a countable set. 
\item <5-> In this case, 
\EQ{
A_X(x) = \cap_{t \in T}A_{X_t}(x_t) = \cap_{s \in S}A_{X_s}(x_s) = A_{X_S}(x_S) \in \sF.    
}
\end{itemize}
\end{block}
\end{frame}

\begin{frame}{Measurability}
\begin{Example}[Bernoulli sequence] 
\begin{itemize}
\item <1-> Consider a sample space $\set{H,T}^\N$. 
We define a mapping $X: \Omega \to \set{0,1}^\N$ such that $X_n(\omega) = \SetIn{H}(\omega_n) = \SetIn{\omega_n = H}$. %$X(\omega) = (\SetIn{S}(\omega_1), \SetIn{S}(\omega_2), \dots)$. 
\item <2-> The map $X$ is an $\sF$-measurable random sequence, if each $X_n: \Omega \to \set{0,1}$ is a bi-variate $\sF$-measurable random variable on the probability space $(\Omega,\sF,P)$. 
\item <3-> Therefore, the event space $\sF$ must contain the event space generated by sequence of events $E\in\sF^\N$ defined by $E_n \triangleq \set{\omega \in \Omega: X_n(\omega)=1} = \set{\omega \in \Omega: \omega_n = H} \in \sF$ for all $n\in \N$. 
\item <4-> That is, 
\EQ{\sigma(X) = \sigma(E) = \sigma(\set{E_n: n \in \N}).}  
\end{itemize}
\end{Example}
\end{frame}


\begin{frame}{Distribution}
\begin{Definition}<1->
For a random process $X: \Omega \to \sX^T$ defined on the probability space $(\Omega,\sF, P)$, 
we define a \textbf{finite dimensional distribution} $F_{X_S}: \R^S \to [0,1]$ for a finite $S \subseteq T$ 
by 
\EQ{
F_{X_S}(x_S) \triangleq P(A_{X_S}(x_S)),\quad x_S \in \R^S. 
}
\end{Definition}
\begin{Example}<2->
Consider a probability space $(\Omega,\sF, P)$ defined by the sample space $\Omega = \set{H,T}^\N$, the event space $\sF \triangleq \sigma(E)$ where $E_n = \set{\omega \in \Omega: \omega_n = H}$ for $n\in\N$, 
and the probability measure $P: \sF \to [0,1]$ defined by 
\EQ{
P(\cap_{i \in F}E_i) = p^{\abs{F}},\text{ for all finite }F \subseteq \N.
}
\end{Example}
\end{frame}

\begin{frame}{Distribution}
\begin{Example}[Continued]
\begin{itemize}
\item <1-> Let $X: \Omega\to\set{0,1}^\N$ defined as $X_n(\omega) = \Ind{E_n}(\omega)$ for all outcomes $\omega \in \Omega$ and $n \in \N$. 
\item <2-> For this random sequence, we can obtain the finite dimensional distribution $F_{X_S}: \R^S \to [0,1]$ for any finite $S \subseteq T$ and $x \in \R^S$ in terms of $I_0(x) \triangleq \set{i \in S: x_i < 0}$ and $I_1(x) \triangleq \set{i \in S: x_i \in [0,1)}$, 
as 
\item <3-> \EQN{
\label{eqn:FDDBernoulli}
F_{X_S}(x) = 
\begin{cases}
1,& I_0(x) \cup I_1(x) = \emptyset,\\
(1-p)^{\abs{I_1(x)}}, & I_0(x) = \emptyset, I_1(x) \neq \emptyset,\\
0, & I_0(x) \neq \emptyset.
\end{cases}
}
\end{itemize}
\end{Example}
\end{frame}

\begin{frame}{Joint Distribution}
\begin{itemize}
\item <1-> We are interested in identifying the joint distribution $F: \R^T \to [0,1]$. 
To this end, for any $x \in \R^T$, we need to know
\EQ{
F_X(x) \triangleq P\left(\displaystyle {\bigcap_{t \in T}\{\omega \in \Omega: X_t(\omega) \le x_t\}}\right) 
= P(\bigcap_{t \in T}X_t^{-1}(-\infty, x_t]) 
= P(A_X(x)).
}
\item <2-> We can verify that $A_X(x) \in \sF$ for $x \in \R^T$ such that $\set{t \in T: x_t < \infty}$ is countable. 
\item <3-> For a simple independent process with countably infinite $T$, any function of the above form would be zero if $x_t$ is finite for all $t \in T$. 
\item <4-> Therefore, we only look at the values of $F_X(x)$ for $x \in \R^T$ where $\set{t \in T: x_t < \infty}$ is finite. 
\item <5-> Set of all finite dimensional distributions of the stochastic process $X: \Omega \to \sX^T$ characterizes its distribution completely. 
\end{itemize}
\end{frame}

\begin{frame}{Joint Distirbution}
\begin{Example}
\begin{itemize}
\item <1-> Consider a probability space $(\Omega,\sF, P)$ defined by the sample space $\Omega = \set{H,T}^\N$ and the event space $\sF \triangleq \sigma(E)$ where $E_n = \set{\omega \in \Omega: \omega_n = H}$ for all $n\in\N$. 
\item <2-> Let $X: \Omega\to\set{0,1}^\N$ defined as $X_n(\omega) = \Ind{E_n}(\omega)$ for all outcomes $\omega \in \Omega$ and $n \in \N$. 
\item <3-> For this random sequence, if we are given the finite dimensional distribution $F_{X_S}: \R^S \to [0,1]$ for any finite $S \subseteq T$ and $x \in \R^S$ in terms of sets $ I_0(x) \triangleq \set{i \in S: x_i < 0}$ and $I_1(x) \triangleq \set{i \in S: x_i \in [0,1)}$, 
as defined in Eq.~\eqref{eqn:FDDBernoulli}. 
\item <4-> Then, we can find the probability measure $P: \sF \to [0,1]$ is given by 
\EQ{
P(\cap_{i \in F}E_i) = p^{\abs{F}},\text{ for all finite }F \subseteq \N.
}
\end{itemize}
\end{Example}
\end{frame}

\begin{frame}{Independence}
\begin{Definition} <1->
A random process is \textbf{independent} if the collection of event spaces $(\sigma(X_t): t \in T)$ is independent. 
That is, for all $x_S \in \R^S$, we have 
\EQ{
F_{X_S}(x_S) = P(\cap_{s \in S}\set{X_s \le  x_s}) 
= \prod_{s \in S}P\set{X_s \le x_s} 
= \prod_{s \in S}F_{X_s}(x_s). 
}
\end{Definition}

\begin{block}{Remark}<2->
Independence of a random process is equivalent to factorization of any finite dimensional distribution function into product of individual marginal distribution functions.
\end{block}
\end{frame}

\begin{frame}{Independence}
\begin{Example}
\begin{itemize}
\item <1-> Consider a probability space $(\Omega,\sF, P)$ defined by the sample space $\Omega = \set{H,T}^\N$, the event space $\sF \triangleq \sigma(E)$ where $E_n = \set{\omega \in \Omega: \omega_n = H}$ for all $n\in\N$, 
and the probability measure $P: \sF \to [0,1]$ defined by 
\EQ{
P(\cap_{i \in F}E_i) = p^{\abs{F}},\text{ for all finite }F \subseteq \N.
}
\item <2-> Then, we observe that the random sequence $X: \Omega\to\set{0,1}^\N$ defined by $X_n(\omega) \triangleq \Ind{E_n}(\omega)$ for all outcomes $\omega \in \Omega$ and $n \in \N$, is independent.
\end{itemize}
\end{Example}
\end{frame}

\begin{frame}{Independence}
\begin{Definition}<1->
Two stochastic processes $X: \Omega \to \sX^{T_1},Y:\Omega\to \sY^{T_2}$ are \textbf{independent}, if the corresponding event spaces $\sigma(X), \sigma(Y)$ are independent. 
\end{Definition}
\begin{Remark}
\begin{itemize}
    \item <2-> For any $x \in \R^{S_1}, y \in \R^{S_2}$ for finite $S_1 \subseteq T_1, S_2\subseteq T_2$, 
    the events $A_{S_1}(x) \triangleq \cap_{s\in S_1}X_s^{-1}(-\infty, x_s]$ and $B_{S_2}(y) \triangleq \cap_{s \in S_2}Y_s^{-1}(-\infty, y_s]$ are independent. 
    \item <3-> The joint finite dimensional distribution of $X$ and $Y$ factorizes, and 
    \EQ{
    P(A_{S_1}(x)\cap B_{S_2}(y)) 
    = P(A_{S_1}(x))P(B_{S_2}(y)) 
    = F_{X_{S_1}}(x)F_{Y_{S_2}}(y),\\
    \quad x \in \R^{S_1}, y \in \R^{S_2}.
    }
\end{itemize}
\end{Remark}
\end{frame}
\end{document}
